\section{Iteration and the thin spanning tree}\label{sec:iteration}

We now turn the sparsification step into a uniform bound.
The estimate keeps track of the tree-packing parameter rather than
the number of rounds.
The bounded product of losses uses the summable-error principle in
the spectral-halving argument of Harvey and Olver
\cite[Appendix~F, Claims~F.3--F.4]{HarveyOlver2014}.

\begin{proof}[Proof of Theorem~\ref{thm:main}]
Choose an absolute integer \(r_*\) at least the cutoff in
Proposition~\ref{prop:step}. For every \(r\ge r_*\), that proposition
supplies a positive integer \(r'\) satisfying
\[
 \frac5{16}\le\frac{r'}r\le\frac12,
 \qquad
 \frac{1/2+3\tau}{r'/r}\le1+15r^{-1/16}.
\]
These are exactly the estimates in \eqref{eq:step-ratio-bound}.

First let \(H_0\) be any graph on at least two vertices with
\(r_0\ge1\) packed spanning trees. As long as \(r_i\ge r_*\), apply
Proposition~\ref{prop:step} to obtain a subgraph \(H_{i+1}\)
with at least
\[
 r_{i+1}=
 \left\lfloor\frac{(1-4r_i^{-1/16})r_i}{2}\right\rfloor
\]
packed spanning trees. Stop at the first \(j\) with \(r_j<r_*\).
This happens after finitely many steps because
\(r_{i+1}\le r_i/2\). Every \(r_i\) is positive.

If \(j\ge1\), reading the halving inequality backwards gives
\[
 r_i\ge2^{j-1-i}r_{j-1}
 \ge2^{j-1-i}r_*
 \qquad(0\le i<j).
\]
Consequently
\[
 \sum_{i=0}^{j-1}r_i^{-1/16}
 \le\frac{r_*^{-1/16}}{1-2^{-1/16}}.
\]
Set
\begin{equation}\label{eq:constant-C2}
 C_{\mathrm{it}}=\exp\left(
 \frac{15r_*^{-1/16}}{1-2^{-1/16}}
 \right)\ge1.
\end{equation}
For every cut \(S\), multiplication of the step inequalities and
\eqref{eq:step-ratio-bound} gives
\[
 \begin{aligned}
 |\delta_{H_j}(S)|
 &\le\frac{r_j}{r_0}
   \prod_{i=0}^{j-1}(1+15r_i^{-1/16})
   |\delta_{H_0}(S)|\\
 &\le C_{\mathrm{it}}\frac{r_j}{r_0}|\delta_{H_0}(S)|.
 \end{aligned}
\]
For \(j=0\) the same inequality holds with an empty product.
All these estimates concern the same sequence of subgraphs and
hold for every cut simultaneously.

Since \(r_j\ge1\), the final graph \(H_j\) is connected.
Any spanning tree \(T\) of \(H_j\) satisfies
\begin{equation}\label{eq:packing-thinness}
 |\delta_T(S)|\le\frac{C_{\mathrm{it}}r_*}{r_0}|\delta_{H_0}(S)|
 \qquad(\varnothing\ne S\subsetneq V).
\end{equation}
This includes \(r_0<r_*\): there are then zero rounds, and any
spanning tree works because \(C_{\mathrm{it}}r_*/r_0\ge1\).

Finally take \(H_0=G\). If \(k\ge2\), summing the cut lower bounds
over a partition into \(d\ge2\) parts shows that
\[
 |E_G(\mathcal P)|\ge\frac{kd}{2}
 \ge\lfloor k/2\rfloor(d-1).
\]
Theorem~\ref{thm:packing} supplies \(\lfloor k/2\rfloor\) packed
trees. If \(k=1\), connectedness supplies one tree.
Thus in all cases we may choose
\[
 r_0=\max\{1,\lfloor k/2\rfloor\}\ge k/3.
\]
Substitution in \eqref{eq:packing-thinness} proves the theorem with
the universal constant
\[
 C=3C_{\mathrm{it}}r_*.
\]
\end{proof}
